\begin{abstract}
In regression models with subject-level random effects, the likelihood
identifies the sum of a subject-constant fixed effect and the matching
combination of random effects, but not its split. Under the
non-centered parameterization that Hamiltonian Monte Carlo (HMC) software
uses by default, only the prior locates the split, and the resulting ridge
can slow sampling of the reported coefficients. We
characterize every random-effect direction the fixed effects can absorb,
and the condition a survival submodel adds in joint
longitudinal-survival models. One fixed orthogonal rotation of the
standardized random effects gathers these directions into a small block of
sampling coordinates. It is exact for any random-effect covariance, and a dense mass-matrix block over this block and the fixed
effects largely whitens the ridge. The effective sample size per second of
the intercept and time slope rose 14- and 21-fold in simulation
(geometric means), and that
of the intercept and subject-constant coefficients 17- to 28-fold on two
joint-model datasets. For the association parameter of a joint model, an
ideal two-block sampler's lag-1 autocorrelation equals the share of its
posterior variance explained by the other parameters, which on real data
was at least 0.85. Supplementary materials for this article are available
online.
\end{abstract}

\noindent%
{\it Keywords:} Identifiability; Longitudinal and time-to-event data;
Mixed models; No-U-Turn sampler; Non-centered parameterization;
Preconditioning.
\vfill
